\begin{abcliste}
\abc With $\ssc{E}{kin} = \frac{1}{2}m_\alpha v^2$ (the energy in joules):
\begin{align}
 v &= \vF \\
   &= \sqrt{\frac{2\times\Ea\times\nce}{\ma}} \\
   &= \v \approx \result{\vP{0}{2}}
\end{align}
about \vcpP{0}{2} of the speed of light.

\abc
\begin{align}
 \lambda &= \frac{h}{m_\alpha v} = \laF \\
   &= \frac{\nch}{\sqrt{2\times\ma\times\Ea\times\nce}} \\
   &= \la \approx \result{\laP{-15}{2}}
\end{align}
About the size of the nucleus (\SI{7}{\femto\m}), so its wave nature matters there.

\abc Classically it $\result{\text{could never leave the nucleus}}$: the wall is about \EbO\ high, far above \EaO, so it would bounce back for ever. Quantum mechanically, each time it hits the wall it has a tiny chance to tunnel through: about $10^{21}$ hits per second, and a chance that makes the half-life 1600~years.

\abc At \EpO\ the alpha particle meets the wall higher up, where the barrier is $\result{\text{lower above its energy and thinner}}$. The transmission grows exponentially, and the half-life shrinks from 1600~years to \SI{0.3}{\micro\s}. (The alpha particle does not break through the wall: \EpO\ is still far below its top.) Gamow explained alpha decay this way in 1928, the first success of tunnelling.
\end{abcliste}
